
3.10
Soit \(A\) une \(K\)-\(\lambda\)-algèbre augmentée. On définit sur \(A\) une filtration d'anneau de la manière suivante : pour tout \(n\ge0\), on désigne par \(\Fil^n A\) l'idéal engendré par les produits :
\begin{equation}\tag{3.10.1}
  \gamma^{i_1}(x_1)\ldots\gamma^{i_k}(x_k),
  \qquad \text{avec } i_1+\ldots+i_k\ge n,\ \text{et } \forall i,\ \epsilon(x_i)=0
\end{equation}
\((\epsilon:A\to K\) étant l'augmentation).

Supposons que \(A\) soit engendré en tant que groupe abélien par une famille d'éléments \((x_\alpha)_{\alpha\in I}\)~; alors \(\Fil^n A\) est engendré comme \(K\)-module par les produits :
\begin{equation}\tag{3.10.2}
  \gamma^{i_1}(x_{\alpha_1}-\epsilon(x_{\alpha_1}))\ldots
  \gamma^{i_k}(x_{\alpha_k}-\epsilon(x_{\alpha_k})),
  \qquad \text{avec } i_1+\ldots+i_k\ge n,\ \text{et } \alpha_i\in I
\end{equation}
En effet, si \(x\) est un élément de \(A\) d'augmentation nulle et si on écrit \(x=\sum_i x_{\alpha_i}\), on a \(\sum_i \epsilon(x_{\alpha_i})=0\), donc \(x=\sum_i x_{\alpha_i}-\epsilon(x_{\alpha_i})\)~; par ailleurs, si \(a\in A\), on peut écrire \(a=(a-\epsilon(a))+\epsilon(a)=\gamma^1(a-\epsilon(a))+\epsilon(a)\)~; en utilisant les relations (3.4.1), on voit donc que les éléments de \(\Fil^n A\) sont combinaisons linéaires à coefficients dans \(K\) des éléments (3.10.2)~; comme ceux-ci sont dans \(\Fil^n A\), on obtient le résultat voulu.

La filtration ainsi définie est une filtration d'anneau~; elle est décroissante, et on a :
\[
  \Fil^0A=A~;
  \qquad \Fil^1A=\Ker\epsilon.
\]
De plus, elle est fonctorielle par rapport à \(A\) : tout \(\lambda\)-homomorphisme de \(K\)-\(\lambda\)-algèbres augmentées est compatible aux filtrations. Quand nous considérerons un \(\lambda\)-anneau augmenté comme muni d'une filtration, il s'agira, sauf mention expresse du contraire, de la filtration ci-dessus.

